% Options for packages loaded elsewhere
% Options for packages loaded elsewhere
\PassOptionsToPackage{unicode}{hyperref}
\PassOptionsToPackage{hyphens}{url}
\PassOptionsToPackage{dvipsnames,svgnames,x11names}{xcolor}
%
\documentclass[
  letterpaper,
]{report}
\usepackage{xcolor}
\usepackage{amsmath,amssymb}
\setcounter{secnumdepth}{5}
\usepackage{iftex}
\ifPDFTeX
  \usepackage[T1]{fontenc}
  \usepackage[utf8]{inputenc}
  \usepackage{textcomp} % provide euro and other symbols
\else % if luatex or xetex
  \usepackage{unicode-math} % this also loads fontspec
  \defaultfontfeatures{Scale=MatchLowercase}
  \defaultfontfeatures[\rmfamily]{Ligatures=TeX,Scale=1}
\fi
\usepackage{lmodern}
\ifPDFTeX\else
  % xetex/luatex font selection
\fi
% Use upquote if available, for straight quotes in verbatim environments
\IfFileExists{upquote.sty}{\usepackage{upquote}}{}
\IfFileExists{microtype.sty}{% use microtype if available
  \usepackage[]{microtype}
  \UseMicrotypeSet[protrusion]{basicmath} % disable protrusion for tt fonts
}{}
\makeatletter
\@ifundefined{KOMAClassName}{% if non-KOMA class
  \IfFileExists{parskip.sty}{%
    \usepackage{parskip}
  }{% else
    \setlength{\parindent}{0pt}
    \setlength{\parskip}{6pt plus 2pt minus 1pt}}
}{% if KOMA class
  \KOMAoptions{parskip=half}}
\makeatother
% Make \paragraph and \subparagraph free-standing
\makeatletter
\ifx\paragraph\undefined\else
  \let\oldparagraph\paragraph
  \renewcommand{\paragraph}{
    \@ifstar
      \xxxParagraphStar
      \xxxParagraphNoStar
  }
  \newcommand{\xxxParagraphStar}[1]{\oldparagraph*{#1}\mbox{}}
  \newcommand{\xxxParagraphNoStar}[1]{\oldparagraph{#1}\mbox{}}
\fi
\ifx\subparagraph\undefined\else
  \let\oldsubparagraph\subparagraph
  \renewcommand{\subparagraph}{
    \@ifstar
      \xxxSubParagraphStar
      \xxxSubParagraphNoStar
  }
  \newcommand{\xxxSubParagraphStar}[1]{\oldsubparagraph*{#1}\mbox{}}
  \newcommand{\xxxSubParagraphNoStar}[1]{\oldsubparagraph{#1}\mbox{}}
\fi
\makeatother

\usepackage{color}
\usepackage{fancyvrb}
\newcommand{\VerbBar}{|}
\newcommand{\VERB}{\Verb[commandchars=\\\{\}]}
\DefineVerbatimEnvironment{Highlighting}{Verbatim}{commandchars=\\\{\}}
% Add ',fontsize=\small' for more characters per line
\usepackage{framed}
\definecolor{shadecolor}{RGB}{241,243,245}
\newenvironment{Shaded}{\begin{snugshade}}{\end{snugshade}}
\newcommand{\AlertTok}[1]{\textcolor[rgb]{0.68,0.00,0.00}{#1}}
\newcommand{\AnnotationTok}[1]{\textcolor[rgb]{0.37,0.37,0.37}{#1}}
\newcommand{\AttributeTok}[1]{\textcolor[rgb]{0.40,0.45,0.13}{#1}}
\newcommand{\BaseNTok}[1]{\textcolor[rgb]{0.68,0.00,0.00}{#1}}
\newcommand{\BuiltInTok}[1]{\textcolor[rgb]{0.00,0.23,0.31}{#1}}
\newcommand{\CharTok}[1]{\textcolor[rgb]{0.13,0.47,0.30}{#1}}
\newcommand{\CommentTok}[1]{\textcolor[rgb]{0.37,0.37,0.37}{#1}}
\newcommand{\CommentVarTok}[1]{\textcolor[rgb]{0.37,0.37,0.37}{\textit{#1}}}
\newcommand{\ConstantTok}[1]{\textcolor[rgb]{0.56,0.35,0.01}{#1}}
\newcommand{\ControlFlowTok}[1]{\textcolor[rgb]{0.00,0.23,0.31}{\textbf{#1}}}
\newcommand{\DataTypeTok}[1]{\textcolor[rgb]{0.68,0.00,0.00}{#1}}
\newcommand{\DecValTok}[1]{\textcolor[rgb]{0.68,0.00,0.00}{#1}}
\newcommand{\DocumentationTok}[1]{\textcolor[rgb]{0.37,0.37,0.37}{\textit{#1}}}
\newcommand{\ErrorTok}[1]{\textcolor[rgb]{0.68,0.00,0.00}{#1}}
\newcommand{\ExtensionTok}[1]{\textcolor[rgb]{0.00,0.23,0.31}{#1}}
\newcommand{\FloatTok}[1]{\textcolor[rgb]{0.68,0.00,0.00}{#1}}
\newcommand{\FunctionTok}[1]{\textcolor[rgb]{0.28,0.35,0.67}{#1}}
\newcommand{\ImportTok}[1]{\textcolor[rgb]{0.00,0.46,0.62}{#1}}
\newcommand{\InformationTok}[1]{\textcolor[rgb]{0.37,0.37,0.37}{#1}}
\newcommand{\KeywordTok}[1]{\textcolor[rgb]{0.00,0.23,0.31}{\textbf{#1}}}
\newcommand{\NormalTok}[1]{\textcolor[rgb]{0.00,0.23,0.31}{#1}}
\newcommand{\OperatorTok}[1]{\textcolor[rgb]{0.37,0.37,0.37}{#1}}
\newcommand{\OtherTok}[1]{\textcolor[rgb]{0.00,0.23,0.31}{#1}}
\newcommand{\PreprocessorTok}[1]{\textcolor[rgb]{0.68,0.00,0.00}{#1}}
\newcommand{\RegionMarkerTok}[1]{\textcolor[rgb]{0.00,0.23,0.31}{#1}}
\newcommand{\SpecialCharTok}[1]{\textcolor[rgb]{0.37,0.37,0.37}{#1}}
\newcommand{\SpecialStringTok}[1]{\textcolor[rgb]{0.13,0.47,0.30}{#1}}
\newcommand{\StringTok}[1]{\textcolor[rgb]{0.13,0.47,0.30}{#1}}
\newcommand{\VariableTok}[1]{\textcolor[rgb]{0.07,0.07,0.07}{#1}}
\newcommand{\VerbatimStringTok}[1]{\textcolor[rgb]{0.13,0.47,0.30}{#1}}
\newcommand{\WarningTok}[1]{\textcolor[rgb]{0.37,0.37,0.37}{\textit{#1}}}

\usepackage{longtable,booktabs,array}
\usepackage{calc} % for calculating minipage widths
% Correct order of tables after \paragraph or \subparagraph
\usepackage{etoolbox}
\makeatletter
\patchcmd\longtable{\par}{\if@noskipsec\mbox{}\fi\par}{}{}
\makeatother
% Allow footnotes in longtable head/foot
\IfFileExists{footnotehyper.sty}{\usepackage{footnotehyper}}{\usepackage{footnote}}
\makesavenoteenv{longtable}
\usepackage{graphicx}
\makeatletter
\newsavebox\pandoc@box
\newcommand*\pandocbounded[1]{% scales image to fit in text height/width
  \sbox\pandoc@box{#1}%
  \Gscale@div\@tempa{\textheight}{\dimexpr\ht\pandoc@box+\dp\pandoc@box\relax}%
  \Gscale@div\@tempb{\linewidth}{\wd\pandoc@box}%
  \ifdim\@tempb\p@<\@tempa\p@\let\@tempa\@tempb\fi% select the smaller of both
  \ifdim\@tempa\p@<\p@\scalebox{\@tempa}{\usebox\pandoc@box}%
  \else\usebox{\pandoc@box}%
  \fi%
}
% Set default figure placement to htbp
\def\fps@figure{htbp}
\makeatother





\setlength{\emergencystretch}{3em} % prevent overfull lines

\providecommand{\tightlist}{%
  \setlength{\itemsep}{0pt}\setlength{\parskip}{0pt}}



 


\makeatletter
\@ifpackageloaded{bookmark}{}{\usepackage{bookmark}}
\makeatother
\makeatletter
\@ifpackageloaded{caption}{}{\usepackage{caption}}
\AtBeginDocument{%
\ifdefined\contentsname
  \renewcommand*\contentsname{Table of contents}
\else
  \newcommand\contentsname{Table of contents}
\fi
\ifdefined\listfigurename
  \renewcommand*\listfigurename{List of Figures}
\else
  \newcommand\listfigurename{List of Figures}
\fi
\ifdefined\listtablename
  \renewcommand*\listtablename{List of Tables}
\else
  \newcommand\listtablename{List of Tables}
\fi
\ifdefined\figurename
  \renewcommand*\figurename{Figure}
\else
  \newcommand\figurename{Figure}
\fi
\ifdefined\tablename
  \renewcommand*\tablename{Table}
\else
  \newcommand\tablename{Table}
\fi
}
\@ifpackageloaded{float}{}{\usepackage{float}}
\floatstyle{ruled}
\@ifundefined{c@chapter}{\newfloat{codelisting}{h}{lop}}{\newfloat{codelisting}{h}{lop}[chapter]}
\floatname{codelisting}{Listing}
\newcommand*\listoflistings{\listof{codelisting}{List of Listings}}
\makeatother
\makeatletter
\makeatother
\makeatletter
\@ifpackageloaded{caption}{}{\usepackage{caption}}
\@ifpackageloaded{subcaption}{}{\usepackage{subcaption}}
\makeatother
\usepackage{bookmark}
\IfFileExists{xurl.sty}{\usepackage{xurl}}{} % add URL line breaks if available
\urlstyle{same}
\hypersetup{
  pdftitle={Pureflow Architecture Guide},
  colorlinks=true,
  linkcolor={blue},
  filecolor={Maroon},
  citecolor={Blue},
  urlcolor={Blue},
  pdfcreator={LaTeX via pandoc}}


\title{Pureflow Architecture Guide}
\author{}
\date{}
\begin{document}
\maketitle

\renewcommand*\contentsname{Table of contents}
{
\hypersetup{linkcolor=}
\setcounter{tocdepth}{2}
\tableofcontents
}

\bookmarksetup{startatroot}

\chapter{Pureflow Architecture Guide}\label{pureflow-architecture-guide}

\section{Audience}\label{audience}

This guide is for engineers who need a system-level understanding of
Pureflow.

\section{Scope}\label{scope}

\begin{itemize}
\tightlist
\item
  High-level architecture and subsystem responsibilities.
\item
  Runtime flow and execution model.
\item
  Package and contract boundaries.
\item
  Diagram-backed explanations for key pathways.
\end{itemize}

\section{Relationship to Other Docs}\label{relationship-to-other-docs}

\begin{itemize}
\tightlist
\item
  API reference is generated separately via \texttt{cargo\ doc}.
\item
  Retrieval-oriented facts live in \texttt{docs/llm/}.
\end{itemize}

\bookmarksetup{startatroot}

\chapter{System Overview}\label{system-overview}

\section{What Pureflow Is}\label{what-pureflow-is}

Pureflow is an experimental Flow-Based Programming engine in Rust. It
validates workflow documents, executes node graphs through bounded
channels, and emits machine-facing metadata and run summaries.

At a high level, Pureflow separates:

\begin{itemize}
\tightlist
\item
  static graph and format validation,
\item
  runtime execution and orchestration,
\item
  contract/capability policy,
\item
  introspection and CLI surfaces.
\end{itemize}

\section{Layered Architecture}\label{layered-architecture}

\begin{Shaded}
\begin{Highlighting}[]
\NormalTok{flowchart TD}
\NormalTok{    WF[Workflow Documents\textbackslash{}nJSON/TOML/YAML] {-}{-}\textgreater{} WFF[pureflow{-}workflow{-}format\textbackslash{}nRaw parse + version checks]}
\NormalTok{    WFF {-}{-}\textgreater{} W[pureflow{-}workflow\textbackslash{}nGraph structure validation]}

\NormalTok{    W {-}{-}\textgreater{} C[pureflow{-}contract\textbackslash{}nContracts + schema compatibility]}
\NormalTok{    W {-}{-}\textgreater{} I[pureflow{-}introspection\textbackslash{}nPure projections]}
\NormalTok{    C {-}{-}\textgreater{} I}
\NormalTok{    CAP[pureflow{-}core capability model] {-}{-}\textgreater{} C}

\NormalTok{    C {-}{-}\textgreater{} E[pureflow{-}engine\textbackslash{}nWorkflow orchestration]}
\NormalTok{    CORE[pureflow{-}core\textbackslash{}nPorts, metadata, errors, traits] {-}{-}\textgreater{} E}
\NormalTok{    R[pureflow{-}runtime\textbackslash{}nasupersync adapter] {-}{-}\textgreater{} E}
\NormalTok{    WASM[pureflow{-}wasm\textbackslash{}nWasmtime batch adapter] {-}{-}\textgreater{} E}

\NormalTok{    E {-}{-}\textgreater{} CLI[pureflow{-}cli\textbackslash{}nvalidate/inspect/explain/run]}
\NormalTok{    I {-}{-}\textgreater{} CLI}
\end{Highlighting}
\end{Shaded}

\section{Crate Responsibilities}\label{crate-responsibilities}

\begin{longtable}[]{@{}
  >{\raggedright\arraybackslash}p{(\linewidth - 4\tabcolsep) * \real{0.3333}}
  >{\raggedright\arraybackslash}p{(\linewidth - 4\tabcolsep) * \real{0.3333}}
  >{\raggedright\arraybackslash}p{(\linewidth - 4\tabcolsep) * \real{0.3333}}@{}}
\toprule\noalign{}
\begin{minipage}[b]{\linewidth}\raggedright
Crate
\end{minipage} & \begin{minipage}[b]{\linewidth}\raggedright
Primary responsibility
\end{minipage} & \begin{minipage}[b]{\linewidth}\raggedright
Used by
\end{minipage} \\
\midrule\noalign{}
\endhead
\bottomrule\noalign{}
\endlastfoot
\texttt{pureflow-types} & Validated identifier primitives
(\texttt{WorkflowId}, \texttt{NodeId}, \texttt{PortId}) & all crates \\
\texttt{pureflow-workflow-format} & External format decode + versioning
(\texttt{pureflow\_version}) & CLI + loaders \\
\texttt{pureflow-workflow} & Static topology model + structural
validation & contract/introspection/engine \\
\texttt{pureflow-core} & Runtime-facing traits, ports, metadata, errors,
capabilities & runtime/engine/cli/wasm \\
\texttt{pureflow-contract} & Node contract model + validation against
workflow/capabilities & engine/introspection/cli \\
\texttt{pureflow-introspection} & Read-only workflow+contract+capability
projections & CLI inspect/explain \\
\texttt{pureflow-runtime} & \texttt{asupersync} runtime wrapper + node
observer boundary & engine \\
\texttt{pureflow-engine} & Workflow orchestration, registry scheduling,
summary/policy & CLI \\
\texttt{pureflow-wasm} & Wasmtime Component Model batch execution adapter
& engine/CLI \\
\texttt{pureflow-cli} & User and automation entrypoint & humans/CI \\
\texttt{pureflow-test-kit} & Builders and helpers for tests/examples &
tests/docs \\
\end{longtable}

\section{Runtime Data Flow}\label{runtime-data-flow}

\begin{Shaded}
\begin{Highlighting}[]
\NormalTok{sequenceDiagram}
\NormalTok{    participant User}
\NormalTok{    participant CLI as pureflow{-}cli}
\NormalTok{    participant Format as workflow{-}format}
\NormalTok{    participant Graph as workflow/contract}
\NormalTok{    participant Engine as pureflow{-}engine}
\NormalTok{    participant Runtime as pureflow{-}runtime}
\NormalTok{    participant Node as Native or WASM node}
\NormalTok{    participant Sink as Metadata JSONL}

\NormalTok{    User{-}\textgreater{}\textgreater{}CLI: run workflow + metadata path}
\NormalTok{    CLI{-}\textgreater{}\textgreater{}Format: parse workflow document}
\NormalTok{    Format{-}\textgreater{}\textgreater{}Graph: build validated workflow definition}
\NormalTok{    CLI{-}\textgreater{}\textgreater{}Graph: validate contracts + capabilities}
\NormalTok{    CLI{-}\textgreater{}\textgreater{}Engine: run\_workflow\_with\_registry...}
\NormalTok{    Engine{-}\textgreater{}\textgreater{}Runtime: execute nodes via observers}
\NormalTok{    Runtime{-}\textgreater{}\textgreater{}Node: NodeExecutor::run / BatchExecutor::invoke}
\NormalTok{    Node{-}{-}\textgreater{}\textgreater{}Runtime: packets or error}
\NormalTok{    Runtime{-}\textgreater{}\textgreater{}Sink: lifecycle/message/queue/error records}
\NormalTok{    Engine{-}{-}\textgreater{}\textgreater{}CLI: WorkflowRunSummary + terminal state}
\NormalTok{    CLI{-}{-}\textgreater{}\textgreater{}User: text or {-}{-}json run summary}
\end{Highlighting}
\end{Shaded}

\section{Metadata Surfaces and Usage}\label{metadata-surfaces-and-usage}

Pureflow emits two machine-facing outputs:

\begin{itemize}
\tightlist
\item
  metadata JSONL stream (event records),
\item
  \texttt{pureflow\ run\ -\/-json} terminal summary.
\end{itemize}

\subsection{Metadata JSONL record
families}\label{metadata-jsonl-record-families}

\begin{longtable}[]{@{}
  >{\raggedright\arraybackslash}p{(\linewidth - 6\tabcolsep) * \real{0.2500}}
  >{\raggedright\arraybackslash}p{(\linewidth - 6\tabcolsep) * \real{0.2500}}
  >{\raggedright\arraybackslash}p{(\linewidth - 6\tabcolsep) * \real{0.2500}}
  >{\raggedright\arraybackslash}p{(\linewidth - 6\tabcolsep) * \real{0.2500}}@{}}
\toprule\noalign{}
\begin{minipage}[b]{\linewidth}\raggedright
\texttt{record\_type}
\end{minipage} & \begin{minipage}[b]{\linewidth}\raggedright
Produced by
\end{minipage} & \begin{minipage}[b]{\linewidth}\raggedright
Purpose
\end{minipage} & \begin{minipage}[b]{\linewidth}\raggedright
Common consumers
\end{minipage} \\
\midrule\noalign{}
\endhead
\bottomrule\noalign{}
\endlastfoot
\texttt{execution\_context} & runtime observer boundary & Tie records to
workflow/node/attempt/cancellation state & debugging, correlation \\
\texttt{lifecycle} & runtime + engine & Node
scheduling/start/completion/failure/cancel transitions & health
dashboards, run trace \\
\texttt{message} & port boundary instrumentation & Packet movement
(\texttt{enqueued}/\texttt{dequeued}/\texttt{dropped}) without payload
bytes & flow auditing \\
\texttt{queue\_pressure} & port reserve/recv boundaries & Capacity and
queue state diagnostics (\texttt{reserve\_full},
\texttt{receive\_empty}, etc.) & backpressure analysis \\
\texttt{error} & node/workflow failure paths & Stable error taxonomy
(\texttt{CDT-*}) with optional diagnostics & alerting, retry policy \\
\texttt{external\_effect} & node integration code & Explicit
tool/service/database/API effect observations & governance, audit
trails \\
\end{longtable}

\subsection{Run summary JSON fields}\label{run-summary-json-fields}

\begin{longtable}[]{@{}
  >{\raggedright\arraybackslash}p{(\linewidth - 4\tabcolsep) * \real{0.3333}}
  >{\raggedright\arraybackslash}p{(\linewidth - 4\tabcolsep) * \real{0.3333}}
  >{\raggedright\arraybackslash}p{(\linewidth - 4\tabcolsep) * \real{0.3333}}@{}}
\toprule\noalign{}
\begin{minipage}[b]{\linewidth}\raggedright
Field
\end{minipage} & \begin{minipage}[b]{\linewidth}\raggedright
Meaning
\end{minipage} & \begin{minipage}[b]{\linewidth}\raggedright
Notes
\end{minipage} \\
\midrule\noalign{}
\endhead
\bottomrule\noalign{}
\endlastfoot
\texttt{status} & terminal state (\texttt{completed}, \texttt{failed},
\texttt{cancelled}) & primary machine decision field \\
\texttt{error} & terminal error object or \texttt{null} & mirrors stable
Pureflow error shape \\
\texttt{workflow} & workflow identity and shape counts & quick context
for automation \\
\texttt{metadata} & output metadata path and record totals & ties
summary to JSONL artifact \\
\texttt{summary} & node scheduling/completion/failure counts & execution
outcome snapshot \\
\end{longtable}

\section{Important Examples}\label{important-examples}

\subsection{CLI usage}\label{cli-usage}

\begin{Shaded}
\begin{Highlighting}[]
\CommentTok{\# Validate workflow structure and format}
\ExtensionTok{cargo}\NormalTok{ run }\AttributeTok{{-}p}\NormalTok{ pureflow{-}cli }\AttributeTok{{-}{-}}\NormalTok{ validate examples/native{-}linear{-}etl.workflow.json}

\CommentTok{\# Inspect topology/contracts/capabilities}
\ExtensionTok{cargo}\NormalTok{ run }\AttributeTok{{-}p}\NormalTok{ pureflow{-}cli }\AttributeTok{{-}{-}}\NormalTok{ inspect examples/native{-}linear{-}etl.workflow.json}

\CommentTok{\# Run and emit machine{-}facing summary}
\ExtensionTok{cargo}\NormalTok{ run }\AttributeTok{{-}p}\NormalTok{ pureflow{-}cli }\AttributeTok{{-}{-}}\NormalTok{ run }\AttributeTok{{-}{-}json} \DataTypeTok{\textbackslash{}}
\NormalTok{  examples/native{-}linear{-}etl.workflow.json }\DataTypeTok{\textbackslash{}}
\NormalTok{  /tmp/pureflow{-}native{-}linear.metadata.jsonl}
\end{Highlighting}
\end{Shaded}

\subsection{Workflow document shape}\label{workflow-document-shape}

\begin{Shaded}
\begin{Highlighting}[]
\FunctionTok{\{}
  \DataTypeTok{"pureflow\_version"}\FunctionTok{:} \StringTok{"1"}\FunctionTok{,}
  \DataTypeTok{"id"}\FunctionTok{:} \StringTok{"native{-}linear{-}etl"}\FunctionTok{,}
  \DataTypeTok{"nodes"}\FunctionTok{:} \OtherTok{[}
    \FunctionTok{\{} \DataTypeTok{"id"}\FunctionTok{:} \StringTok{"source"}\FunctionTok{,} \DataTypeTok{"inputs"}\FunctionTok{:} \OtherTok{[]}\FunctionTok{,} \DataTypeTok{"outputs"}\FunctionTok{:} \OtherTok{[}\StringTok{"rows"}\OtherTok{]} \FunctionTok{\}}\OtherTok{,}
    \FunctionTok{\{} \DataTypeTok{"id"}\FunctionTok{:} \StringTok{"transform"}\FunctionTok{,} \DataTypeTok{"inputs"}\FunctionTok{:} \OtherTok{[}\StringTok{"rows"}\OtherTok{]}\FunctionTok{,} \DataTypeTok{"outputs"}\FunctionTok{:} \OtherTok{[}\StringTok{"cleaned"}\OtherTok{]} \FunctionTok{\}}\OtherTok{,}
    \FunctionTok{\{} \DataTypeTok{"id"}\FunctionTok{:} \StringTok{"sink"}\FunctionTok{,} \DataTypeTok{"inputs"}\FunctionTok{:} \OtherTok{[}\StringTok{"cleaned"}\OtherTok{]}\FunctionTok{,} \DataTypeTok{"outputs"}\FunctionTok{:} \OtherTok{[]} \FunctionTok{\}}
  \OtherTok{]}\FunctionTok{,}
  \DataTypeTok{"edges"}\FunctionTok{:} \OtherTok{[}
    \FunctionTok{\{}
      \DataTypeTok{"source"}\FunctionTok{:} \FunctionTok{\{} \DataTypeTok{"node"}\FunctionTok{:} \StringTok{"source"}\FunctionTok{,} \DataTypeTok{"port"}\FunctionTok{:} \StringTok{"rows"} \FunctionTok{\},}
      \DataTypeTok{"target"}\FunctionTok{:} \FunctionTok{\{} \DataTypeTok{"node"}\FunctionTok{:} \StringTok{"transform"}\FunctionTok{,} \DataTypeTok{"port"}\FunctionTok{:} \StringTok{"rows"} \FunctionTok{\},}
      \DataTypeTok{"capacity"}\FunctionTok{:} \DecValTok{2}
    \FunctionTok{\}}
  \OtherTok{]}
\FunctionTok{\}}
\end{Highlighting}
\end{Shaded}

\subsection{Contract/capability alignment example
(Rust)}\label{contractcapability-alignment-example-rust}

\begin{Shaded}
\begin{Highlighting}[]
\KeywordTok{use} \PreprocessorTok{pureflow\_contract::}\OperatorTok{\{}
\NormalTok{    Determinism}\OperatorTok{,}\NormalTok{ ExecutionMode}\OperatorTok{,}\NormalTok{ NodeContract}\OperatorTok{,}\NormalTok{ PortContract}\OperatorTok{,}\NormalTok{ SchemaRef}\OperatorTok{,}
\NormalTok{    validate\_workflow\_contracts}\OperatorTok{,}
\OperatorTok{\};}
\KeywordTok{use} \PreprocessorTok{pureflow\_core::}\OperatorTok{\{}
\NormalTok{    RetryDisposition}\OperatorTok{,}
    \PreprocessorTok{capability::}\OperatorTok{\{}\NormalTok{NodeCapabilities}\OperatorTok{,}\NormalTok{ PortCapability}\OperatorTok{,}\NormalTok{ PortCapabilityDirection}\OperatorTok{\},}
\OperatorTok{\};}

\CommentTok{// Build contracts and capabilities for each workflow node, then validate:}
\NormalTok{validate\_workflow\_contracts(}\OperatorTok{\&}\NormalTok{workflow}\OperatorTok{,} \OperatorTok{\&}\NormalTok{contracts}\OperatorTok{,} \OperatorTok{\&}\NormalTok{capabilities)}\OperatorTok{?;}
\end{Highlighting}
\end{Shaded}

This is a key authoring boundary for both humans and future LLM
assistants: workflow topology, contracts, and capabilities must match
before execution.

\section{Current Boundary Decisions}\label{current-boundary-decisions}

\begin{itemize}
\tightlist
\item
  \texttt{pureflow-workflow} validates structural graph honesty only.
\item
  \texttt{pureflow-contract} validates contract/capability alignment plus
  schema compatibility across connected ports.
\item
  \texttt{pureflow-core} owns runtime interfaces and metadata shape.
\item
  \texttt{pureflow-runtime} is an adapter over \texttt{asupersync}, not
  the public model.
\item
  \texttt{pureflow-wasm} isolates Wasmtime and WIT/component ABI details
  from core.
\end{itemize}

\section{Areas Needing Clarification (For Follow-up
Review)}\label{areas-needing-clarification-for-follow-up-review}

These are the points where code and docs suggest open questions or
evolving policy. I recommend focused follow-up review on each.

\begin{enumerate}
\def\labelenumi{\arabic{enumi}.}
\item
  Feedback-loop runtime behavior beyond startup mode. The engine exposes
  \texttt{CycleRunPolicy::AllowFeedbackLoops} and watchdog policies, but
  expected production semantics for complex cyclic workloads are not yet
  clearly documented.
\item
  Process execution mode roadmap. \texttt{ExecutionMode::Process} exists
  as a contract enum variant, but adapter, enforcement, and metadata
  expectations are future-facing.
\item
  External effect emission conventions. The metadata model supports
  \texttt{external\_effect} records, but there is still room to
  standardize operation/target naming guidance for cross-team
  consistency.
\item
  Queue-pressure interpretation guidance. Queue-pressure kinds are
  well-defined, but threshold recommendations and SLO-style
  interpretation rules are not yet documented in one place.
\item
  WASM capability boundary evolution. Current WASM validation rejects
  declared effects in import-free mode. If host imports are introduced
  later, capability enforcement rules and migration expectations will
  need a dedicated update.
\end{enumerate}

\bookmarksetup{startatroot}

\chapter{Runtime Model}\label{runtime-model}

\section{Intent}\label{intent}

Explain how execution proceeds from input through node processing to
output.

\section{Sections to Fill}\label{sections-to-fill}

\begin{itemize}
\tightlist
\item
  Runtime lifecycle and orchestration.
\item
  Execution phases and state transitions.
\item
  Error handling and observability touchpoints.
\end{itemize}

\bookmarksetup{startatroot}

\chapter{Package and Contract
Boundaries}\label{package-and-contract-boundaries}

\section{Intent}\label{intent-1}

Define crate-level roles and key contract surfaces.

\section{Sections to Fill}\label{sections-to-fill-1}

\begin{itemize}
\tightlist
\item
  Crate map and responsibility boundaries.
\item
  Public contract surfaces and invariants.
\item
  Compatibility rules and versioning expectations.
\end{itemize}

\bookmarksetup{startatroot}

\chapter{Diagrams and Flows}\label{diagrams-and-flows}

\section{Intent}\label{intent-2}

Track architecture diagrams and flow descriptions tied to narrative
chapters.

\section{Sections to Fill}\label{sections-to-fill-2}

\begin{itemize}
\tightlist
\item
  Diagram inventory and source location.
\item
  Diagram generation notes for \texttt{frankenmermaid}.
\item
  Canonical flows to represent first.
\end{itemize}




\end{document}
